\section{Scelte implementative}
\label{sec:implementazione}

\subsection{Mappa locale e conoscenza parziale}

Ogni agente mantiene una propria rappresentazione interna dell'ambiente: una matrice $25 \times 25$ inizializzata con la \textbf{griglia reale} caricata dal file JSON di input. Tutti gli agenti conoscono quindi fin dall'inizio la posizione di muri, corridoi, magazzini e porte direzionali: la struttura della mappa non è mai parzialmente osservabile.

La vera informazione nascosta sono le \textbf{posizioni degli oggetti}: questi non fanno parte della griglia ma sono gestiti separatamente dall'environment. Un oggetto viene reso visibile a un agente solo quando \texttt{observe()} lo rilevisce all'interno del raggio visivo (\texttt{reveal\_object\_at}), dopodiché viene aggiunto al registro locale \texttt{known\_objects}.

Per tracciare cosa rimane da esplorare, ogni agente mantiene un insieme \texttt{\_unobserved}: l'insieme di tutte le celle non-muro che non sono ancora state attraversate entro il raggio visivo. Gli Scout si dirigono verso queste celle non per scoprire la topologia (già nota), ma per \textbf{passarci abbastanza vicino da rivelare eventuali oggetti nascosti}. La fase di esplorazione è dunque in realtà una fase di \emph{scansione} del territorio alla ricerca di oggetti.

\subsection{Comunicazione peer-to-peer}

Quando due agenti si trovano entro il raggio di comunicazione (\emph{comm\_radius}), eseguono un \emph{merge} bidirezionale delle rispettive conoscenze. Poiché la struttura della mappa è già nota a tutti, il merge riguarda essenzialmente due informazioni:

\begin{itemize}
  \item Il \textbf{registro degli oggetti noti} (\texttt{known\_objects}): gli oggetti osservati da uno Scout vengono trasmessi all'altro agente, rendendoli immediatamente disponibili per il fetch;
  \item L'\textbf{insieme delle celle non ancora visitate} (\texttt{\_unobserved}): viene sostituito con l'intersezione dei due insiemi, così le celle già scansionate da uno non vengono ri-visitate dall'altro, evitando ridondanza.
\end{itemize}

Vengono anche condivisi i \texttt{claimed\_frontier}, le celle-obiettivo già prese in carico, così gli Scout si coordinano per non dirigersi verso la stessa zona.

Questa comunicazione locale -- senza infrastruttura centralizzata -- è il meccanismo attraverso cui l'intelligenza collettiva emerge: un Scout che ha scansionato il quadrante nord-est trasferisce immediatamente quella conoscenza a qualsiasi agente che gli si avvicini.

\subsection{Pathfinding con A*}

Il modulo \texttt{pathfinding.py} implementa l'algoritmo A* con euristica di distanza di Manhattan. L'algoritmo riceve una funzione \texttt{is\_walkable\_fn} che opera sulla mappa locale dell'agente anziché sulla mappa reale: questo garantisce che la pianificazione sia coerente con la conoscenza parziale disponibile.

La scelta di A* rispetto ad algoritmi più semplici (BFS, greedy) è motivata dalla necessità di:
\begin{itemize}
  \item Gestire rotte ottimali in presenza di ostacoli complessi e passaggi stretti;
  \item Rispettare i vincoli direzionali di ingresso/uscita dei magazzini;
  \item Operare efficientemente su una griglia di dimensione $25 \times 25$ (625 celle), dove BFS sarebbe comunque accettabile, ma A* offre un margine migliore in scenari con molti muri.
\end{itemize}

Quando A* non trova un percorso (es. la mappa locale non ha ancora rivelato il collegamento verso la destinazione), l'agente si sposta casualmente verso la destinazione in attesa di scoprire nuove celle.

\subsection{Strategia di esplorazione degli Scout}

Gli Scout adottano una strategia di \emph{coverage-based scanning} con repulsione inter-agente. L'obiettivo è selezionare la cella non ancora visitata (\texttt{\_unobserved}) che:
\begin{enumerate}
  \item È raggiungibile tramite A*;
  \item Massimizza la distanza media dagli altri Scout attivi.
\end{enumerate}

Questo garantisce una \textbf{scansione uniforme e non ridondante} dell'area. Senza repulsione, più Scout tenderebbero a convergere sulle stesse zone, lasciando ampie aree non scansionate e quindi con oggetti non ancora scoperti.

All'avvio, gli Scout vengono assegnati a quadranti distinti della griglia (angolo superiore-sinistro, superiore-destro, ecc.) mediante un target iniziale fisso, accelerando la diversificazione geografica nei primi tick.

\subsection{Gestione della batteria ed emergenza}

Ogni agente tiene traccia della batteria residua e calcola ad ogni tick la distanza stimata fino al magazzino più vicino accessibile. Se:
\[
\text{batteria} \leq \text{distanza\_magazzino} + \text{margine\_sicurezza}
\]
l'agente entra in stato \texttt{EMERGENCY} e abbandona il task corrente per tornare al magazzino, indipendentemente da ciò che stava facendo (incluso trasportare un oggetto).

Il margine di sicurezza è differenziato per ruolo: i Collector hanno un margine maggiore (20 passi) rispetto agli Scout (10 passi), poiché il rischio di perdere un oggetto in transito è più costoso del rischio di perdere uno Scout esploratore.

Gli agenti che muoiono lontano da un magazzino (batteria zero senza essere rientrati) vengono marcati come \texttt{DEAD} e non emettono più azioni, ma le loro ultime posizioni rimangono visibili nella visualizzazione per debug.

\subsection{Visualizzazione real-time}

La visualizzazione Pygame mostra la mappa dal punto di vista dell'\textbf{unione delle mappe locali} di tutti gli agenti, cioè ciò che il collettivo conosce nell'istante corrente. Le celle ancora \texttt{UNKNOWN} appaiono in grigio scuro. Questo permette di osservare l'esplorazione collettiva emergere tick per tick.

Ogni agente è rappresentato con un colore specifico (verde per Scout, arancione per Collector, viola per Relay, grigio per Dead), e le celle di ingresso/uscita dei magazzini mostrano frecce direzionali. Il pannello laterale mostra le metriche in tempo reale: tick corrente, batteria di ciascun agente, oggetti consegnati e rimanenti.

\subsection{Logging e analisi post-simulazione}

Ogni esecuzione produce:
\begin{itemize}
  \item \textbf{\texttt{simulation\_log.json}}: snapshot completo ad ogni tick (posizioni agenti, stati, oggetti residui);
  \item \textbf{\texttt{results.json}}: metriche scalari aggregate (ticks totali, energia media per agente, energia per oggetto, ticks per oggetto);
  \item \textbf{\texttt{heatmap.png}}: mappa di calore della frequenza di visita per cella, aggregata su tutti gli agenti e tutti i tick.
\end{itemize}

Il modulo \texttt{run\_all.py} esegue in sequenza tutti e 6 i scenari (2 mappe $\times$ 3 configurazioni) in modalità headless e produce al termine quattro grafici comparativi aggregati.
